\subsection{现有方法的证明检查：从模型到可证伪命题}
\subsection{参数、机制与预测}
任何模型都应把参数分为偏好参数、技术参数和制度参数。参数不是“越多越好”：每增加一个自由参数，就必须说明其识别来源以及对预测的边际贡献。一个可证伪命题至少包含对象、参数变化、结果变量和方向判断。

\subsection{局部与全局}
局部比较静态围绕基准点展开，依赖导数和非奇异雅可比矩阵；全局结论则要求凹性、单调性或格理论等结构条件。二者不能混用。若论文只证明了局部结论，应明确写出“在给定邻域内”。

\subsection{研究设计检查表}
\begin{enumerate}
\item 可行集是否非空、闭且有界？\item 目标函数是否连续？\item 最优解是否唯一？\item 参数是否可由数据识别？\item 反事实是否超出支持集？\item 结论是否依赖某个不可检验的偏好限制？
\end{enumerate}
\subsection{讨论题}
请分别以价格补贴、配额和信息披露为例，写出可行集变化、最优集合变化以及福利判定所需的额外条件。每个案例均应区分正向预测与规范评价。
